% Licence, CC-BY
\documentclass[compress, ucs]{beamer}\usepackage[]{graphicx}\usepackage[]{color}
%% maxwidth is the original width if it is less than linewidth
%% otherwise use linewidth (to make sure the graphics do not exceed the margin)
\makeatletter
\def\maxwidth{ %
  \ifdim\Gin@nat@width>\linewidth
    \linewidth
  \else
    \Gin@nat@width
  \fi
}
\makeatother

\definecolor{fgcolor}{rgb}{0.345, 0.345, 0.345}
\newcommand{\hlnum}[1]{\textcolor[rgb]{0.686,0.059,0.569}{#1}}%
\newcommand{\hlstr}[1]{\textcolor[rgb]{0.192,0.494,0.8}{#1}}%
\newcommand{\hlcom}[1]{\textcolor[rgb]{0.678,0.584,0.686}{\textit{#1}}}%
\newcommand{\hlopt}[1]{\textcolor[rgb]{0,0,0}{#1}}%
\newcommand{\hlstd}[1]{\textcolor[rgb]{0.345,0.345,0.345}{#1}}%
\newcommand{\hlkwa}[1]{\textcolor[rgb]{0.161,0.373,0.58}{\textbf{#1}}}%
\newcommand{\hlkwb}[1]{\textcolor[rgb]{0.69,0.353,0.396}{#1}}%
\newcommand{\hlkwc}[1]{\textcolor[rgb]{0.333,0.667,0.333}{#1}}%
\newcommand{\hlkwd}[1]{\textcolor[rgb]{0.737,0.353,0.396}{\textbf{#1}}}%

\usepackage{framed}
\makeatletter
\newenvironment{kframe}{%
 \def\at@end@of@kframe{}%
 \ifinner\ifhmode%
  \def\at@end@of@kframe{\end{minipage}}%
  \begin{minipage}{\columnwidth}%
 \fi\fi%
 \def\FrameCommand##1{\hskip\@totalleftmargin \hskip-\fboxsep
 \colorbox{shadecolor}{##1}\hskip-\fboxsep
     % There is no \\@totalrightmargin, so:
     \hskip-\linewidth \hskip-\@totalleftmargin \hskip\columnwidth}%
 \MakeFramed {\advance\hsize-\width
   \@totalleftmargin\z@ \linewidth\hsize
   \@setminipage}}%
 {\par\unskip\endMakeFramed%
 \at@end@of@kframe}
\makeatother

\definecolor{shadecolor}{rgb}{.97, .97, .97}
\definecolor{messagecolor}{rgb}{0, 0, 0}
\definecolor{warningcolor}{rgb}{1, 0, 1}
\definecolor{errorcolor}{rgb}{1, 0, 0}
\newenvironment{knitrout}{}{} % an empty environment to be redefined in TeX

\usepackage{alltt}
\usepackage[T2A]{fontenc}
\usepackage[utf8x]{inputenc}
\usepackage[english, russian]{babel}
\usepackage{amssymb}
\usepackage{url}
% Breaking long urls on hyphens
\makeatletter
\g@addto@macro{\UrlBreaks}{\UrlOrds}
\makeatother
\usepackage{pscyr}
\usepackage{graphicx}
\usetheme{Warsaw}
\usecolortheme{beaver}
\usefonttheme{professionalfonts}
\beamertemplatenavigationsymbolsempty
\setbeamertemplate{headline}{}
\expandafter\def\expandafter\insertshorttitle\expandafter{%
\insertshorttitle\hfill%
\insertframenumber\,/\,\inserttotalframenumber}

\makeatletter
\definecolor{beaverred}{RGB}{142, 19, 19}
\setbeamertemplate{enumerate items}{\color{beaverred}$\blacklozenge$}
\setbeamertemplate{itemize items}{\color{beaverred}$\blacklozenge$}
\setbeamertemplate{section in toc}[square]
\setbeamertemplate{subsection in toc}[square]
\makeatother
% --------------------------------------------------------------
% Setting up some knitr options

% --------------------------------------------------------------
\IfFileExists{upquote.sty}{\usepackage{upquote}}{}
\begin{document}
\title[dplyr/tidyr/magrittr]{\normalsize Современный взгляд на предобработку данных: \\
dplyr, tidyr и magrittr}
\author{Антон Антонов}
\date{28 марта 2015}
%\institute{\textbf {SPb R User Group}}
\titlegraphic{\includegraphics[width=2cm]{logo}}
%\email{tonytonov@gmail.com}
\maketitle
% --------------------------------------------------------------
\begin{frame}[fragile]{Цикл жизни данных в R}
\begin{center}
\includegraphics[width=\linewidth]{cycle}
\end{center}
\pause
\begin{itemize}
\item tidyr 0.2 (Hadley Wickham, RStudio)
\item dplyr 0.4 (Hadley Wickham, Roman Francois, RStudio)
\item magrittr 1.5 (Stefan Milton Bache, Hadley Wickham)
\end{itemize}
\end{frame}
% --------------------------------------------------------------
\begin{frame}[fragile]{magrittr::\%>\%}
При необходимости последовательно вызвать несколько функций подряд есть следующие возможности.
\pause
\begin{itemize}
\item Вложенные вызовы (nested calls):
\begin{knitrout}\scriptsize
\definecolor{shadecolor}{rgb}{0.882, 0.882, 0.882}\color{fgcolor}\begin{kframe}
\begin{alltt}
\hlkwd{arrange}\hlstd{(}
  \hlkwd{summarize}\hlstd{(}
    \hlkwd{filter}\hlstd{(data, variable} \hlopt{==} \hlnum{42}\hlstd{),}
    \hlkwc{Total} \hlstd{=} \hlkwd{sum}\hlstd{(variable)}
  \hlstd{),}
  \hlkwd{desc}\hlstd{(Total)}
\hlstd{)}
\end{alltt}
\end{kframe}
\end{knitrout}
\pause
\item Создание промежуточных переменных:
\begin{knitrout}\scriptsize
\definecolor{shadecolor}{rgb}{0.882, 0.882, 0.882}\color{fgcolor}\begin{kframe}
\begin{alltt}
\hlstd{a} \hlkwb{<-} \hlkwd{filter}\hlstd{(data, variable} \hlopt{==} \hlnum{42}\hlstd{)}
\hlstd{b} \hlkwb{<-} \hlkwd{summarise}\hlstd{(a,} \hlkwc{Total} \hlstd{=} \hlkwd{sum}\hlstd{(variable))}
\hlstd{c} \hlkwb{<-} \hlkwd{arrange}\hlstd{(b,} \hlkwd{desc}\hlstd{(Total))}
\end{alltt}
\end{kframe}
\end{knitrout}
\pause
\item Конвейеры (pipes) при помощи оператора \%>\%:
\begin{knitrout}\scriptsize
\definecolor{shadecolor}{rgb}{0.882, 0.882, 0.882}\color{fgcolor}\begin{kframe}
\begin{alltt}
\hlstd{data} \hlopt{%>%}
  \hlkwd{filter}\hlstd{(variable} \hlopt{==} \hlnum{42}\hlstd{)} \hlopt{%>%}
  \hlkwd{summarise}\hlstd{(}\hlkwc{Total} \hlstd{=} \hlkwd{sum}\hlstd{(variable))} \hlopt{%>%}
  \hlkwd{arrange}\hlstd{(}\hlkwd{desc}\hlstd{(Total))}
\end{alltt}
\end{kframe}
\end{knitrout}
\end{itemize}
\end{frame}
% --------------------------------------------------------------
\begin{frame}[fragile]{magrittr::\%>\%}
Следующий кусок кода содержит вложенные вызовы, формулу с точкой и анонимную функцию.
\begin{knitrout}\scriptsize
\definecolor{shadecolor}{rgb}{0.882, 0.882, 0.882}\color{fgcolor}\begin{kframe}
\begin{alltt}
\hlstd{car_data} \hlkwb{<-}
  \hlkwd{subset}\hlstd{(}
    \hlkwd{transform}\hlstd{(}
      \hlkwd{aggregate}\hlstd{(}
        \hlstd{.} \hlopt{~} \hlstd{cyl,}
        \hlkwc{data} \hlstd{= mtcars[mtcars}\hlopt{$}\hlstd{hp} \hlopt{>} \hlnum{100}\hlstd{, ],}
        \hlkwc{FUN} \hlstd{=} \hlkwa{function}\hlstd{(}\hlkwc{x}\hlstd{)} \hlkwd{round}\hlstd{(}\hlkwd{mean}\hlstd{(x),} \hlnum{2}\hlstd{)}
      \hlstd{),}
      \hlkwc{kpl} \hlstd{= mpg} \hlopt{*} \hlnum{0.4251}
    \hlstd{),}
    \hlkwc{select} \hlstd{=} \hlkwd{c}\hlstd{(cyl, hp, kpl)}
  \hlstd{)}
\hlkwd{print}\hlstd{(car_data)}
\end{alltt}
\begin{verbatim}
##   cyl     hp       kpl
## 1   4 111.00 11.010090
## 2   6 122.29  8.391474
## 3   8 209.21  6.419010
\end{verbatim}
\end{kframe}
\end{knitrout}
\end{frame}
% --------------------------------------------------------------
\begin{frame}[fragile]{magrittr::\%>\%}
Используем оператор, работающий по двум основным правилам:
\begin{knitrout}\scriptsize
\definecolor{shadecolor}{rgb}{0.882, 0.882, 0.882}\color{fgcolor}\begin{kframe}
\begin{alltt}
\hlstd{x} \hlopt{%>%} \hlkwd{f}\hlstd{(y)}            \hlcom{# same as f(x, y)}
\hlstd{x} \hlopt{%>%} \hlkwd{g}\hlstd{(y,} \hlkwc{param} \hlstd{= .)} \hlcom{# same as g(y, param = x)}
\end{alltt}
\end{kframe}
\end{knitrout}
\pause
Тот же код в виде конвейера:
\begin{knitrout}\scriptsize
\definecolor{shadecolor}{rgb}{0.882, 0.882, 0.882}\color{fgcolor}\begin{kframe}
\begin{alltt}
\hlkwd{library}\hlstd{(magrittr)}
\hlstd{car_data} \hlkwb{<-}
  \hlstd{mtcars} \hlopt{%>%}
  \hlkwd{subset}\hlstd{(hp} \hlopt{>} \hlnum{100}\hlstd{)} \hlopt{%>%}
  \hlkwd{aggregate}\hlstd{(.} \hlopt{~} \hlstd{cyl,} \hlkwc{data} \hlstd{= .,}
            \hlkwc{FUN} \hlstd{=} \hlkwa{function}\hlstd{(}\hlkwc{x}\hlstd{)} \hlkwd{round}\hlstd{(}\hlkwd{mean}\hlstd{(x),} \hlnum{2}\hlstd{))} \hlopt{%>%}
  \hlkwd{transform}\hlstd{(}\hlkwc{kpl} \hlstd{= mpg} \hlopt{*} \hlnum{0.4251}\hlstd{)} \hlopt{%>%}
  \hlkwd{subset}\hlstd{(}\hlkwc{select} \hlstd{=} \hlkwd{c}\hlstd{(cyl, hp, kpl))} \hlopt{%>%}
  \hlkwd{print}
\end{alltt}
\begin{verbatim}
##   cyl     hp       kpl
## 1   4 111.00 11.010090
## 2   6 122.29  8.391474
## 3   8 209.21  6.419010
\end{verbatim}
\end{kframe}
\end{knitrout}
\end{frame}
% --------------------------------------------------------------
\begin{frame}[fragile]{magrittr::\%>\%}
Конвейеры можно (но не нужно) использовать практически везде:
\begin{knitrout}\scriptsize
\definecolor{shadecolor}{rgb}{0.882, 0.882, 0.882}\color{fgcolor}\begin{kframe}
\begin{alltt}
\hlstd{car_data} \hlkwb{<-}
  \hlstd{mtcars} \hlopt{%>%}
  \hlkwd{`[`}\hlstd{(.}\hlopt{$}\hlstd{hp} \hlopt{>} \hlnum{100}\hlstd{, )} \hlopt{%>%}
  \hlkwd{aggregate}\hlstd{(.} \hlopt{~} \hlstd{cyl,} \hlkwc{data} \hlstd{= .,}
            \hlkwc{FUN} \hlstd{= .} \hlopt{%>%} \hlstd{mean} \hlopt{%>%} \hlkwd{round}\hlstd{(}\hlnum{2}\hlstd{))} \hlopt{%>%}
  \hlkwd{transform}\hlstd{(}\hlkwc{kpl} \hlstd{= mpg} \hlopt{%>%} \hlkwd{multiply_by}\hlstd{(}\hlnum{0.4251}\hlstd{))} \hlopt{%>%}
  \hlkwd{subset}\hlstd{(}\hlkwc{select} \hlstd{=} \hlkwd{c}\hlstd{(cyl, hp, kpl))} \hlopt{%>%}
  \hlkwd{print}
\end{alltt}
\begin{verbatim}
##   cyl     hp       kpl
## 1   4 111.00 11.010090
## 2   6 122.29  8.391474
## 3   8 209.21  6.419010
\end{verbatim}
\end{kframe}
\end{knitrout}
\end{frame}
% --------------------------------------------------------------
\begin{frame}[fragile]{Tidy data}
Концепция ``tidy data'':
\begin{itemize}
\item каждая колонка (столбец) -- переменная;
\item каждый ряд (строка) -- наблюдение.
\end{itemize}
\pause
\begin{center}
\includegraphics[width=0.5\linewidth]{tidy1}
\end{center}
\par
\pause
Для R такой формат данных наиболее естественен из-за векторизованности многих базовых операций.
\begin{center}
\includegraphics[width=0.3\linewidth]{tidy2}
\end{center}
\end{frame}
% --------------------------------------------------------------
\begin{frame}[fragile]{tidyr}
Возьмём результат некоторого воображаемого эксперимента:

\begin{knitrout}\scriptsize
\definecolor{shadecolor}{rgb}{0.882, 0.882, 0.882}\color{fgcolor}\begin{kframe}
\begin{alltt}
\hlstd{messy}
\end{alltt}
\begin{verbatim}
##   id     group    work.T1   home.T1   work.T2    home.T2
## 1  1 treatment 0.08513597 0.6158293 0.1135090 0.05190332
## 2  2   control 0.22543662 0.4296715 0.5959253 0.26417767
## 3  3 treatment 0.27453052 0.6516557 0.3580500 0.39879073
## 4  4   control 0.27230507 0.5677378 0.4288094 0.83613414
\end{verbatim}
\end{kframe}
\end{knitrout}
\pause
Почему это не ``tidy data''?
\pause
\begin{itemize}
\item Переменные места и времени смешаны;
\item Одной строке соответствует четыре наблюдения.
\end{itemize}
\end{frame}
% --------------------------------------------------------------
\begin{frame}[fragile]{tidyr::gather}
Преобразуем данные из т.н. ``широкого'' формата в ``длинный'' (wide to long):
\begin{knitrout}\scriptsize
\definecolor{shadecolor}{rgb}{0.882, 0.882, 0.882}\color{fgcolor}\begin{kframe}
\begin{alltt}
\hlstd{tidier} \hlkwb{<-} \hlstd{messy} \hlopt{%>%}
  \hlkwd{gather}\hlstd{(key, bpr,} \hlopt{-}\hlstd{id,} \hlopt{-}\hlstd{group)}
\hlstd{tidier} \hlopt{%>%} \hlkwd{head}\hlstd{(}\hlnum{8}\hlstd{)}
\end{alltt}
\begin{verbatim}
##   id     group     key        bpr
## 1  1 treatment work.T1 0.08513597
## 2  2   control work.T1 0.22543662
## 3  3 treatment work.T1 0.27453052
## 4  4   control work.T1 0.27230507
## 5  1 treatment home.T1 0.61582931
## 6  2   control home.T1 0.42967153
## 7  3 treatment home.T1 0.65165567
## 8  4   control home.T1 0.56773775
\end{verbatim}
\end{kframe}
\end{knitrout}
\end{frame}
% --------------------------------------------------------------
\begin{frame}[fragile]{tidyr::spread}
Обратное преобразование (long to wide):
\begin{knitrout}\scriptsize
\definecolor{shadecolor}{rgb}{0.882, 0.882, 0.882}\color{fgcolor}\begin{kframe}
\begin{alltt}
\hlstd{tidier} \hlopt{%>%}
  \hlkwd{spread}\hlstd{(key, bpr)} \hlopt{%>%}
  \hlkwd{print} \hlopt{%>%}
  \hlkwd{all.equal}\hlstd{(messy)}
\end{alltt}
\begin{verbatim}
##   id     group    work.T1   home.T1   work.T2    home.T2
## 1  1 treatment 0.08513597 0.6158293 0.1135090 0.05190332
## 2  2   control 0.22543662 0.4296715 0.5959253 0.26417767
## 3  3 treatment 0.27453052 0.6516557 0.3580500 0.39879073
## 4  4   control 0.27230507 0.5677378 0.4288094 0.83613414
## [1] TRUE
\end{verbatim}
\end{kframe}
\end{knitrout}
\end{frame}
% --------------------------------------------------------------
\begin{frame}[fragile]{tidyr::separate}
Разделяем переменные:
\begin{knitrout}\scriptsize
\definecolor{shadecolor}{rgb}{0.882, 0.882, 0.882}\color{fgcolor}\begin{kframe}
\begin{alltt}
\hlstd{tidy} \hlkwb{<-} \hlstd{tidier} \hlopt{%>%}
  \hlkwd{separate}\hlstd{(key,} \hlkwc{into} \hlstd{=} \hlkwd{c}\hlstd{(}\hlstr{"location"}\hlstd{,} \hlstr{"time"}\hlstd{),} \hlkwc{sep} \hlstd{=} \hlstr{"\textbackslash{}\textbackslash{}."}\hlstd{)}
\hlstd{tidy} \hlopt{%>%} \hlkwd{head}\hlstd{(}\hlnum{8}\hlstd{)}
\end{alltt}
\begin{verbatim}
##   id     group location time        bpr
## 1  1 treatment     work   T1 0.08513597
## 2  2   control     work   T1 0.22543662
## 3  3 treatment     work   T1 0.27453052
## 4  4   control     work   T1 0.27230507
## 5  1 treatment     home   T1 0.61582931
## 6  2   control     home   T1 0.42967153
## 7  3 treatment     home   T1 0.65165567
## 8  4   control     home   T1 0.56773775
\end{verbatim}
\end{kframe}
\end{knitrout}
\end{frame}
% --------------------------------------------------------------
\begin{frame}[fragile]{tidyr::unite}
Склеиваем переменные обратно:
\begin{knitrout}\scriptsize
\definecolor{shadecolor}{rgb}{0.882, 0.882, 0.882}\color{fgcolor}\begin{kframe}
\begin{alltt}
\hlstd{tidy} \hlopt{%>%}
  \hlkwd{unite}\hlstd{(key, location, time,} \hlkwc{sep} \hlstd{=} \hlstr{"."}\hlstd{)} \hlopt{%>%}
  \hlkwd{head}\hlstd{(}\hlnum{8}\hlstd{)}
\end{alltt}
\begin{verbatim}
##   id     group     key        bpr
## 1  1 treatment work.T1 0.08513597
## 2  2   control work.T1 0.22543662
## 3  3 treatment work.T1 0.27453052
## 4  4   control work.T1 0.27230507
## 5  1 treatment home.T1 0.61582931
## 6  2   control home.T1 0.42967153
## 7  3 treatment home.T1 0.65165567
## 8  4   control home.T1 0.56773775
\end{verbatim}
\end{kframe}
\end{knitrout}
\end{frame}
% --------------------------------------------------------------
\begin{frame}[fragile]{dplyr::select}
Выбираем подмножество столбцов (переменных):
\begin{knitrout}\scriptsize
\definecolor{shadecolor}{rgb}{0.882, 0.882, 0.882}\color{fgcolor}\begin{kframe}
\begin{alltt}
\hlstd{tidy} \hlopt{%>%}
  \hlkwd{select}\hlstd{(id, group, location, bpr)} \hlopt{%>%}
  \hlkwd{head}\hlstd{(}\hlnum{8}\hlstd{)}
\end{alltt}
\begin{verbatim}
##   id     group location        bpr
## 1  1 treatment     work 0.08513597
## 2  2   control     work 0.22543662
## 3  3 treatment     work 0.27453052
## 4  4   control     work 0.27230507
## 5  1 treatment     home 0.61582931
## 6  2   control     home 0.42967153
## 7  3 treatment     home 0.65165567
## 8  4   control     home 0.56773775
\end{verbatim}
\end{kframe}
\end{knitrout}
\pause
\begin{knitrout}\scriptsize
\definecolor{shadecolor}{rgb}{0.882, 0.882, 0.882}\color{fgcolor}\begin{kframe}
\begin{alltt}
\hlstd{tidy} \hlopt{%>%} \hlkwd{select}\hlstd{(}\hlnum{1}\hlopt{:}\hlnum{3}\hlstd{,} \hlnum{5}\hlstd{)}
\hlstd{tidy} \hlopt{%>%} \hlkwd{select}\hlstd{(id}\hlopt{:}\hlstd{location,} \hlkwd{starts_with}\hlstd{(}\hlstr{"b"}\hlstd{))}
\hlstd{tidy} \hlopt{%>%} \hlkwd{select}\hlstd{(}\hlopt{-}\hlstd{time)}
\hlstd{tidy} \hlopt{%>%} \hlkwd{select}\hlstd{(}\hlopt{-}\hlkwd{contains}\hlstd{(}\hlstr{"m"}\hlstd{))}
\hlstd{tidy} \hlopt{%>%} \hlkwd{select}\hlstd{(}\hlopt{-}\hlkwd{matches}\hlstd{(}\hlstr{"t.m"}\hlstd{))}
\end{alltt}
\end{kframe}
\end{knitrout}
\end{frame}
% --------------------------------------------------------------
\begin{frame}[fragile]{dplyr::filter}
Выбираем подмножество строк (наблюдений):
\begin{knitrout}\scriptsize
\definecolor{shadecolor}{rgb}{0.882, 0.882, 0.882}\color{fgcolor}\begin{kframe}
\begin{alltt}
\hlstd{tidy} \hlopt{%>%}
  \hlkwd{filter}\hlstd{(group} \hlopt{==} \hlstr{"control"}\hlstd{, bpr} \hlopt{>} \hlnum{0.5}\hlstd{, id} \hlopt{%in%} \hlnum{4}\hlopt{:}\hlnum{5}\hlstd{)}
\end{alltt}
\begin{verbatim}
##   id   group location time       bpr
## 1  4 control     home   T1 0.5677378
## 2  4 control     home   T2 0.8361341
\end{verbatim}
\end{kframe}
\end{knitrout}
\end{frame}
% --------------------------------------------------------------
\begin{frame}[fragile]{dplyr::group\_by, dplyr::summarize}
Агрегация по группам:
\begin{knitrout}\scriptsize
\definecolor{shadecolor}{rgb}{0.882, 0.882, 0.882}\color{fgcolor}\begin{kframe}
\begin{alltt}
\hlstd{tidy} \hlopt{%>%}
  \hlkwd{group_by}\hlstd{(time, location)} \hlopt{%>%}
  \hlkwd{summarize}\hlstd{(}\hlkwc{mean_t} \hlstd{=} \hlkwd{mean}\hlstd{(bpr),} \hlkwc{sd_t} \hlstd{=} \hlkwd{sd}\hlstd{(bpr))}
\end{alltt}
\begin{verbatim}
## Source: local data frame [4 x 4]
## Groups: time
## 
##   time location    mean_t       sd_t
## 1   T1     home 0.5662236 0.09731072
## 2   T1     work 0.2143520 0.08906863
## 3   T2     home 0.3877515 0.33127700
## 4   T2     work 0.3740734 0.20030382
\end{verbatim}
\end{kframe}
\end{knitrout}
\end{frame}
% --------------------------------------------------------------
\begin{frame}[fragile]{dplyr::arrange}
Сортировка (например, в разных направлениях по разным переменным):
\begin{knitrout}\scriptsize
\definecolor{shadecolor}{rgb}{0.882, 0.882, 0.882}\color{fgcolor}\begin{kframe}
\begin{alltt}
\hlstd{tidy} \hlopt{%>%}
  \hlkwd{arrange}\hlstd{(time,} \hlkwd{desc}\hlstd{(bpr))}
\end{alltt}
\begin{verbatim}
##    id     group location time        bpr
## 1   3 treatment     home   T1 0.65165567
## 2   1 treatment     home   T1 0.61582931
## 3   4   control     home   T1 0.56773775
## 4   2   control     home   T1 0.42967153
## 5   3 treatment     work   T1 0.27453052
## 6   4   control     work   T1 0.27230507
## 7   2   control     work   T1 0.22543662
## 8   1 treatment     work   T1 0.08513597
## 9   4   control     home   T2 0.83613414
## 10  2   control     work   T2 0.59592531
## 11  4   control     work   T2 0.42880942
## 12  3 treatment     home   T2 0.39879073
## 13  3 treatment     work   T2 0.35804998
## 14  2   control     home   T2 0.26417767
## 15  1 treatment     work   T2 0.11350898
## 16  1 treatment     home   T2 0.05190332
\end{verbatim}
\end{kframe}
\end{knitrout}
\end{frame}
% --------------------------------------------------------------
\begin{frame}[fragile]{dplyr::mutate}
Создание новых переменных:
\begin{knitrout}\scriptsize
\definecolor{shadecolor}{rgb}{0.882, 0.882, 0.882}\color{fgcolor}\begin{kframe}
\begin{alltt}
\hlstd{tidy} \hlopt{%>%}
  \hlkwd{mutate}\hlstd{(}\hlkwc{perc} \hlstd{= bpr}\hlopt{/}\hlkwd{sum}\hlstd{(bpr),}
         \hlkwc{cperc} \hlstd{=} \hlkwd{cumsum}\hlstd{(perc))} \hlopt{%>%}
  \hlkwd{select}\hlstd{(id, bpr}\hlopt{:}\hlstd{cperc)}
\end{alltt}
\begin{verbatim}
##    id        bpr        perc      cperc
## 1   1 0.08513597 0.013799264 0.01379926
## 2   2 0.22543662 0.036539896 0.05033916
## 3   3 0.27453052 0.044497283 0.09483644
## 4   4 0.27230507 0.044136570 0.13897301
## 5   1 0.61582931 0.099816700 0.23878971
## 6   2 0.42967153 0.069643314 0.30843303
## 7   3 0.65165567 0.105623616 0.41405664
## 8   4 0.56773775 0.092021780 0.50607842
## 9   1 0.11350898 0.018398105 0.52447653
## 10  2 0.59592531 0.096590559 0.62106709
## 11  3 0.35804998 0.058034534 0.67910162
## 12  4 0.42880942 0.069503579 0.74860520
## 13  1 0.05190332 0.008412751 0.75701795
## 14  2 0.26417767 0.042819240 0.79983719
## 15  3 0.39879073 0.064638000 0.86447519
## 16  4 0.83613414 0.135524811 1.00000000
\end{verbatim}
\end{kframe}
\end{knitrout}
\end{frame}
% --------------------------------------------------------------
\begin{frame}[fragile]{SQL joins}
Для иллюстрации join рассмотрим два набора данных:

\begin{knitrout}\scriptsize
\definecolor{shadecolor}{rgb}{0.882, 0.882, 0.882}\color{fgcolor}\begin{kframe}
\begin{alltt}
\hlstd{df1}
\end{alltt}
\begin{verbatim}
##             Name Instrument
## 1   Roger Waters       Bass
## 2  David Gilmour     Guitar
## 3    Syd Barrett       Bass
## 4 Richard Wright  Keyboards
## 5     Nick Mason      Drums
\end{verbatim}
\begin{alltt}
\hlstd{df2}
\end{alltt}
\begin{verbatim}
##   Instrument StrCount
## 1       Bass        4
## 2     Guitar        6
## 3     Violin        4
\end{verbatim}
\end{kframe}
\end{knitrout}
\end{frame}
% --------------------------------------------------------------
\begin{frame}[fragile]{SQL joins}
Всем известные joins:
\begin{knitrout}\scriptsize
\definecolor{shadecolor}{rgb}{0.882, 0.882, 0.882}\color{fgcolor}\begin{kframe}
\begin{alltt}
\hlkwd{left_join}\hlstd{(df1, df2,} \hlstr{"Instrument"}\hlstd{)}
\end{alltt}
\begin{verbatim}
##             Name Instrument StrCount
## 1   Roger Waters       Bass        4
## 2  David Gilmour     Guitar        6
## 3    Syd Barrett       Bass        4
## 4 Richard Wright  Keyboards       NA
## 5     Nick Mason      Drums       NA
\end{verbatim}
\begin{alltt}
\hlkwd{full_join}\hlstd{(df1, df2,} \hlstr{"Instrument"}\hlstd{)}
\end{alltt}
\begin{verbatim}
##             Name Instrument StrCount
## 1   Roger Waters       Bass        4
## 2  David Gilmour     Guitar        6
## 3    Syd Barrett       Bass        4
## 4 Richard Wright  Keyboards       NA
## 5     Nick Mason      Drums       NA
## 6           <NA>     Violin        4
\end{verbatim}
\begin{alltt}
\hlcom{#right_join(df1, df2, "Instrument")}
\hlcom{#inner_join(df1, df2, "Instrument")}
\end{alltt}
\end{kframe}
\end{knitrout}
\end{frame}
% --------------------------------------------------------------
\begin{frame}[fragile]{SQL joins}
Не настолько известные, но тоже joins:
\begin{knitrout}\scriptsize
\definecolor{shadecolor}{rgb}{0.882, 0.882, 0.882}\color{fgcolor}\begin{kframe}
\begin{alltt}
\hlkwd{semi_join}\hlstd{(df1, df2,} \hlstr{"Instrument"}\hlstd{)}
\end{alltt}
\begin{verbatim}
##            Name Instrument
## 1  Roger Waters       Bass
## 2   Syd Barrett       Bass
## 3 David Gilmour     Guitar
\end{verbatim}
\begin{alltt}
\hlkwd{anti_join}\hlstd{(df1, df2,} \hlstr{"Instrument"}\hlstd{)}
\end{alltt}
\begin{verbatim}
##             Name Instrument
## 1     Nick Mason      Drums
## 2 Richard Wright  Keyboards
\end{verbatim}
\end{kframe}
\end{knitrout}
\end{frame}
% --------------------------------------------------------------
\begin{frame}[fragile]{rbenchmarking}
Небольшая демонстрация быстродействия dplyr (трудоёмкие операции написаны на C++).
\begin{knitrout}\scriptsize
\definecolor{shadecolor}{rgb}{0.882, 0.882, 0.882}\color{fgcolor}\begin{kframe}
\begin{alltt}
\hlkwd{library}\hlstd{(rbenchmark)}
\hlkwd{set.seed}\hlstd{(}\hlnum{11}\hlstd{)}
\hlstd{n} \hlkwb{<-} \hlnum{10e6}
\hlstd{m} \hlkwb{<-} \hlnum{100}
\hlstd{d} \hlkwb{<-} \hlkwd{data_frame}\hlstd{(}\hlkwc{x} \hlstd{=} \hlkwd{sample}\hlstd{(m, n,} \hlkwc{replace}\hlstd{=}\hlnum{TRUE}\hlstd{),} \hlkwc{y} \hlstd{=} \hlkwd{runif}\hlstd{(n))}
\hlstd{dm} \hlkwb{<-} \hlkwd{data_frame}\hlstd{(}\hlkwc{x} \hlstd{=} \hlkwd{sample}\hlstd{(m))}
\hlstd{l} \hlkwb{<-} \hlkwd{vector}\hlstd{(}\hlstr{"list"}\hlstd{,} \hlnum{5}\hlstd{)}
\end{alltt}
\end{kframe}
\end{knitrout}
\end{frame}
% --------------------------------------------------------------
\begin{frame}[fragile]{rbenchmarking}
\begin{knitrout}\scriptsize
\definecolor{shadecolor}{rgb}{0.882, 0.882, 0.882}\color{fgcolor}\begin{kframe}
\begin{alltt}
\hlkwd{benchmark}\hlstd{(}
  \hlstd{d[d}\hlopt{$}\hlstd{x}\hlopt{>=}\hlnum{10} \hlopt{&} \hlstd{d}\hlopt{$}\hlstd{x}\hlopt{<}\hlnum{20}\hlstd{,],}
  \hlstd{d} \hlopt{%>%} \hlkwd{filter}\hlstd{(x}\hlopt{>=}\hlnum{10}\hlstd{, x}\hlopt{<}\hlnum{20}\hlstd{),}
  \hlkwc{replications}\hlstd{=}\hlnum{5}\hlstd{,}
  \hlkwc{columns}\hlstd{=}\hlkwd{c}\hlstd{(}\hlstr{"test"}\hlstd{,} \hlstr{"elapsed"}\hlstd{,} \hlstr{"relative"}\hlstd{),}
  \hlkwc{order}\hlstd{=}\hlkwa{NULL}
\hlstd{)} \hlopt{%>%} \hlkwd{print} \hlkwb{->} \hlstd{l[[}\hlnum{1}\hlstd{]]}
\end{alltt}
\begin{verbatim}
##                            test elapsed relative
## 1     d[d$x >= 10 & d$x < 20, ]    3.89    2.016
## 2 d %>% filter(x >= 10, x < 20)    1.93    1.000
\end{verbatim}
\end{kframe}
\end{knitrout}
\end{frame}
% --------------------------------------------------------------
\begin{frame}[fragile]{rbenchmarking}
\begin{knitrout}\scriptsize
\definecolor{shadecolor}{rgb}{0.882, 0.882, 0.882}\color{fgcolor}\begin{kframe}
\begin{alltt}
\hlkwd{benchmark}\hlstd{(}
  \hlstd{d[}\hlkwd{order}\hlstd{(d}\hlopt{$}\hlstd{x), ],}
  \hlstd{d} \hlopt{%>%} \hlkwd{arrange}\hlstd{(x),}
  \hlkwc{replications}\hlstd{=}\hlnum{5}\hlstd{,}
  \hlkwc{columns}\hlstd{=}\hlkwd{c}\hlstd{(}\hlstr{"test"}\hlstd{,} \hlstr{"elapsed"}\hlstd{,} \hlstr{"relative"}\hlstd{),}
  \hlkwc{order}\hlstd{=}\hlkwa{NULL}
\hlstd{)} \hlopt{%>%} \hlkwd{print} \hlkwb{->} \hlstd{l[[}\hlnum{2}\hlstd{]]}
\end{alltt}
\begin{verbatim}
##               test elapsed relative
## 1  d[order(d$x), ]   29.44     1.23
## 2 d %>% arrange(x)   23.94     1.00
\end{verbatim}
\end{kframe}
\end{knitrout}
\end{frame}
% --------------------------------------------------------------
\begin{frame}[fragile]{rbenchmarking}
\begin{knitrout}\scriptsize
\definecolor{shadecolor}{rgb}{0.882, 0.882, 0.882}\color{fgcolor}\begin{kframe}
\begin{alltt}
\hlkwd{benchmark}\hlstd{(}
  \hlstd{d}\hlopt{$}\hlstd{y2} \hlkwb{<-} \hlstd{d}\hlopt{$}\hlstd{y} \hlopt{+} \hlstd{d}\hlopt{$}\hlstd{x}\hlopt{/}\hlnum{2}\hlstd{,}
  \hlstd{d} \hlopt{%>%} \hlkwd{mutate}\hlstd{(}\hlkwc{y2} \hlstd{= y} \hlopt{+} \hlstd{x}\hlopt{/}\hlnum{2}\hlstd{),}
  \hlkwc{replications}\hlstd{=}\hlnum{5}\hlstd{,}
  \hlkwc{columns}\hlstd{=}\hlkwd{c}\hlstd{(}\hlstr{"test"}\hlstd{,} \hlstr{"elapsed"}\hlstd{,} \hlstr{"relative"}\hlstd{),}
  \hlkwc{order}\hlstd{=}\hlkwa{NULL}
\hlstd{)} \hlopt{%>%} \hlkwd{print} \hlkwb{->} \hlstd{l[[}\hlnum{3}\hlstd{]]}
\end{alltt}
\begin{verbatim}
##                         test elapsed relative
## 1        d$y2 <- d$y + d$x/2    0.25     1.00
## 2 d %>% mutate(y2 = y + x/2)    0.28     1.12
\end{verbatim}
\end{kframe}
\end{knitrout}
\end{frame}
% --------------------------------------------------------------
\begin{frame}[fragile]{rbenchmarking}
\begin{knitrout}\scriptsize
\definecolor{shadecolor}{rgb}{0.882, 0.882, 0.882}\color{fgcolor}\begin{kframe}
\begin{alltt}
\hlkwd{benchmark}\hlstd{(}
  \hlkwd{tapply}\hlstd{(d}\hlopt{$}\hlstd{y, d}\hlopt{$}\hlstd{x, mean),}
  \hlstd{d} \hlopt{%>%} \hlkwd{group_by}\hlstd{(x)} \hlopt{%>%} \hlkwd{summarize}\hlstd{(}\hlkwc{ym} \hlstd{=} \hlkwd{mean}\hlstd{(y)),}
  \hlkwc{replications}\hlstd{=}\hlnum{5}\hlstd{,}
  \hlkwc{columns}\hlstd{=}\hlkwd{c}\hlstd{(}\hlstr{"test"}\hlstd{,} \hlstr{"elapsed"}\hlstd{,} \hlstr{"relative"}\hlstd{),}
  \hlkwc{order}\hlstd{=}\hlkwa{NULL}
\hlstd{)} \hlopt{%>%} \hlkwd{print} \hlkwb{->} \hlstd{l[[}\hlnum{4}\hlstd{]]}
\end{alltt}
\begin{verbatim}
##                                            test elapsed
## 1                        tapply(d$y, d$x, mean)    4.20
## 2 d %>% group_by(x) %>% summarize(ym = mean(y))    2.66
##   relative
## 1    1.579
## 2    1.000
\end{verbatim}
\end{kframe}
\end{knitrout}
\end{frame}
% --------------------------------------------------------------
\begin{frame}[fragile]{rbenchmarking}
\begin{knitrout}\scriptsize
\definecolor{shadecolor}{rgb}{0.882, 0.882, 0.882}\color{fgcolor}\begin{kframe}
\begin{alltt}
\hlkwd{benchmark}\hlstd{(}
  \hlkwd{merge}\hlstd{(d, dm,} \hlkwc{by}\hlstd{=}\hlstr{"x"}\hlstd{),}
  \hlstd{d} \hlopt{%>%} \hlkwd{inner_join}\hlstd{(dm,} \hlkwc{by}\hlstd{=}\hlstr{"x"}\hlstd{),}
  \hlkwc{replications}\hlstd{=}\hlnum{1}\hlstd{,}
  \hlkwc{columns}\hlstd{=}\hlkwd{c}\hlstd{(}\hlstr{"test"}\hlstd{,} \hlstr{"elapsed"}\hlstd{,} \hlstr{"relative"}\hlstd{),}
  \hlkwc{order}\hlstd{=}\hlkwa{NULL}
\hlstd{)} \hlopt{%>%} \hlkwd{print} \hlkwb{->} \hlstd{l[[}\hlnum{5}\hlstd{]]}
\end{alltt}
\begin{verbatim}
##                             test elapsed relative
## 1         merge(d, dm, by = "x")   40.33  118.618
## 2 d %>% inner_join(dm, by = "x")    0.34    1.000
\end{verbatim}
\end{kframe}
\end{knitrout}
\end{frame}
% --------------------------------------------------------------
\begin{frame}[fragile]{rbenchmarking}
Сводная таблица быстродействия 
\par
(относительное время исполнения, меньше -- лучше):
\begin{center}
  \begin{tabular}{ c | c | c | c | c | c }
    \hline
          & Filter & Sort & New Variable & Aggregation & Join \\ \hline
    base  & 2.016
          & 1.23
          & 1
          & 1.579
          & 118.618 \\ \hline
    dplyr & 1
          & 1
          & 1.12
          & 1
          & 1 \\
    \hline
  \end{tabular}
\end{center}
\pause
Интересный факт: в Rstudio оператор \%>\% имеет своё сочетание горячих клавиш (Ctrl + Shift + M)! 
\end{frame}
% --------------------------------------------------------------
\begin{frame}[fragile]{magrittr::\%<>\%}
С помощью конвейеров можно создавать не только анонимные функции:
\begin{knitrout}\scriptsize
\definecolor{shadecolor}{rgb}{0.882, 0.882, 0.882}\color{fgcolor}\begin{kframe}
\begin{alltt}
\hlnum{0} \hlopt{%>%} \hlkwd{cos} \hlopt{%>%} \hlkwd{sin}
\end{alltt}
\begin{verbatim}
## [1] 0.841471
\end{verbatim}
\begin{alltt}
\hlstd{f} \hlkwb{<-} \hlstd{.} \hlopt{%>%} \hlkwd{cos} \hlopt{%>%} \hlkwd{sin}
\hlstd{f}\hlstd{(}\hlnum{0}\hlstd{)}
\end{alltt}
\begin{verbatim}
## [1] 0.841471
\end{verbatim}
\begin{alltt}
\hlkwd{sin}\hlstd{(}\hlkwd{cos}\hlstd{(}\hlnum{0}\hlstd{))}
\end{alltt}
\begin{verbatim}
## [1] 0.841471
\end{verbatim}
\end{kframe}
\end{knitrout}
\pause
Есть ещё несколько конвейерных операторов, в том числе присваивающий...
\begin{knitrout}\scriptsize
\definecolor{shadecolor}{rgb}{0.882, 0.882, 0.882}\color{fgcolor}\begin{kframe}
\begin{alltt}
\hlstd{df} \hlkwb{<-} \hlkwd{na.omit}\hlstd{(df)}
\hlstd{df} \hlopt{%<>%} \hlkwd{na.omit}
\end{alltt}
\end{kframe}
\end{knitrout}
\end{frame}
% --------------------------------------------------------------
\begin{frame}[fragile]{magrittr::\%T>\%}
... и Tee operator (производит side-effect):
\begin{knitrout}\scriptsize
\definecolor{shadecolor}{rgb}{0.882, 0.882, 0.882}\color{fgcolor}\begin{kframe}
\begin{alltt}
\hlstd{mtcars} \hlopt{%>%}
  \hlkwd{subset}\hlstd{(cyl} \hlopt{==} \hlnum{4}\hlstd{,} \hlkwd{c}\hlstd{(mpg, hp))} \hlopt{%T>%}
  \hlkwd{plot} \hlopt{%>%}
  \hlkwd{colMeans}
\end{alltt}
\begin{verbatim}
##      mpg       hp 
## 26.66364 82.63636
\end{verbatim}
\end{kframe}
\begin{center}
\includegraphics[width=.75\linewidth]{figure/magrittr-final-3-1} 
\end{center}

\end{knitrout}
\end{frame}
% --------------------------------------------------------------
\begin{frame}{Ссылки}
При подготовке доклада были использованы следующие материалы:
\footnotesize{
\begin{itemize}
\item \url{https://github.com/hadley/tidyr}
\item \url{https://github.com/hadley/dplyr}
\item \url{https://github.com/smbache/magrittr}
\pause
\item \url{http://rpubs.com/bradleyboehmke/data_wrangling}
\item \url{http://www.rstudio.com/wp-content/uploads/2015/02/data-wrangling-cheatsheet.pdf}
\pause
\item \url{http://datascience.la/dplyr-and-a-very-basic-benchmark/}
\item \url{http://zevross.com/blog/2015/01/13/a-new-data-processing-workflow-for-r-dplyr-magrittr-tidyr-ggplot2/}
\end{itemize}
\pause
\begin{center}
\Large{Спасибо!}
\end{center}
}
\end{frame}
% --------------------------------------------------------------
\end{document}
